\section{Use Cases per CTF Manager}\label{use-cases-per-ctf-manager}

\subsection{\texorpdfstring{\texttt{UC-01}: Deploy
dell\textquotesingle infrastruttura di rete per il team
\(k\)}{UC-01: Deploy dell\textquotesingle infrastruttura di rete per il team k}}\label{uc-01-deploy-dellinfrastruttura-di-rete-per-il-team-k}

\begin{itemize}
\tightlist
\item
  \textbf{Attore}: CTF Manager
\item
  \textbf{Obiettivo}: Effettuare il deploy delle interfacce di rete
  necessarie al team \(k\) per eseguire i container e consentire
  l\textquotesingle accesso VPN ai partecipanti.
\item
  \textbf{Precondizione}: BlueAgent e\textquotesingle{} attivo ed
  e\textquotesingle{} raggiungibile via API REST.
\end{itemize}

Flusso delle azioni:

\begin{enumerate}
\def\labelenumi{\arabic{enumi}.}
\tightlist
\item
  Il CTF Manager invia la richiesta
  \passthrough{\lstinline!POST /teams!}, specificando
  l\textquotesingle identificativo \(k\) del team ed il relativo nome.
\item
  BlueAgent verifica se l\textquotesingle infrastruttura per il team
  \(k\) esiste gia:

  \begin{enumerate}
  \def\labelenumii{\arabic{enumii}.}
  \tightlist
  \item
    Se il team e\textquotesingle{} gia registrato, risponde
    \passthrough{\lstinline!HTTP 200 OK!}.
  \item
    Altrimenti, l\textquotesingle infrastruttura di rete viene creata da
    zero:

    \begin{enumerate}
    \def\labelenumiii{\arabic{enumiii}.}
    \tightlist
    \item
      Viene creato il bridge \passthrough{\lstinline!br-team-k!} per il
      team \(k\) con IP per il gateway
      \passthrough{\lstinline!10.k.1.1/24!}.
    \item
      Viene creata l\textquotesingle interfaccia WireGuard
      \passthrough{\lstinline!wg-k!} per il team \(k\). La porta UDP
      e\textquotesingle{} \passthrough{\lstinline!51820+k!},
      l\textquotesingle IP per il gateway e\textquotesingle{}
      \passthrough{\lstinline!10.k.0.1/24!}.
    \item
      Viene registrato il team all\textquotesingle interno del database.
    \item
      BlueAgent risponde con stato
      \passthrough{\lstinline!HTTP 201 Created!} confermando
      l\textquotesingle avvenuta creazione.
    \end{enumerate}
  \end{enumerate}
\end{enumerate}

Se avviene un errore irrimediabile durante la creazione, il bridge e
l\textquotesingle interfaccia WireGuard create vengono rilasciate.

\subsection{\texorpdfstring{\texttt{UC-02}: Deploy della challenge per
il team
\(k\)}{UC-02: Deploy della challenge per il team k}}\label{uc-02-deploy-della-challenge-per-il-team-k}

\begin{itemize}
\tightlist
\item
  \textbf{Attore}: CTF Manager
\item
  \textbf{Obiettivo}: Effettuare il deploy dei container necessari alla
  challenge da allestire per il team \(k\) richiedente.
\item
  \textbf{Precondizione}: \passthrough{\lstinline!UC-01!},
  \passthrough{\lstinline!UC-08!}. E\textquotesingle{} stato eseguito
  con successo il deploy dell\textquotesingle infrastruttura di rete per
  il team \(k\) e le immagini dei container richieste sono presenti nel
  sistema.
\end{itemize}

Flusso delle azioni:

\begin{enumerate}
\def\labelenumi{\arabic{enumi}.}
\tightlist
\item
  Il CTF Manager invia la richiesta
  \passthrough{\lstinline!POST /teams/\{team\_k\}/machines!}
  specificando \passthrough{\lstinline!team\_k!}.
\item
  BlueAgent valida le precondizioni.
\item
  Per ogni container specificato nella configurazione:

  \begin{enumerate}
  \def\labelenumii{\arabic{enumii}.}
  \tightlist
  \item
    Crea l\textquotesingle istanza del container.
  \item
    Collega l\textquotesingle interfaccia virtuale
    \passthrough{\lstinline!eth0!} dell\textquotesingle istanza al
    bridge del team \passthrough{\lstinline!br-team-k!}.
  \item
    Inietta nel filesystem dell\textquotesingle istanza la
    configurazione di rete statica.
  \item
    Avvia l\textquotesingle istanza del container ed applica
    l\textquotesingle indirizzo IP statico.
  \item
    Registra il container nel database.
  \end{enumerate}
\item
  BlueAgent restituisce la risposta con stato
  \passthrough{\lstinline!HTTP 201 Created!}.
\end{enumerate}

Se si verifica un errore durante l\textquotesingle avvio, BlueAgent
arresta ed elimina tutte le istanze e le risorse create fino a quel
momento per la challenge e risponde con un errore.

\subsection{\texorpdfstring{\texttt{UC-03}: Verifica della challenge per
il team
\(k\)}{UC-03: Verifica della challenge per il team k}}\label{uc-03-verifica-della-challenge-per-il-team-k}

\begin{itemize}
\tightlist
\item
  \textbf{Attore}: CTF Manager
\item
  \textbf{Obiettivo}: Verificare l\textquotesingle integrita dei servizi
  e l\textquotesingle efficacia delle patch applicate dal team mediante
  l\textquotesingle esecuzione dei comandi specificati nella
  configurazione della challenge.
\item
  \textbf{Precondizioni}: \passthrough{\lstinline!UC-02!}. I container
  relativi alla challenge specificata per il team \(k\) sono attivi.
\end{itemize}

Flusso delle azioni:

\begin{enumerate}
\def\labelenumi{\arabic{enumi}.}
\tightlist
\item
  Il CTF Manager invia la richiesta
  \passthrough{\lstinline!POST /teams/\{team\_k\}/verificator!}
  trasmettendo la configurazione della challenge da verificare.
\item
  Se una verifica per lo stesso team e\textquotesingle{} gia in corso,
  BlueAgent risponde con \passthrough{\lstinline!HTTP 409 Conflict!}.
\item
  Altrimenti, BlueAgent predispone l\textquotesingle ambiente di
  verifica:

  \begin{enumerate}
  \def\labelenumii{\arabic{enumii}.}
  \tightlist
  \item
    Crea uno switch virtuale denominato
    \passthrough{\lstinline!verify-k!}.
  \item
    Esegue la clonazione di tutti i container della challenge del team
    \(k\).
  \item
    Riconfigura l\textquotesingle interfaccia
    \passthrough{\lstinline!eth0!} dei cloni agganciandola allo switch
    di verifica.
  \item
    Avvia i cloni sandbox e configura al loro interno lo stesso
    indirizzo IP statico di produzione
    \passthrough{\lstinline!10.k.1.\{ip\_suffix\}/24!}.
  \item
    Se specificato nella configurazione, avvia un container di verifica.
  \end{enumerate}
\item
  Esecuzione dei controlli di verifica: viene eseguito un ciclo di
  tentativi finche\textquotesingle{} il comando non ha successo, oppure
  fino allo scadere del timeout.
\item
  Esegue i comandi di verifica all\textquotesingle interno del container
  specificato.
\item
  Dealloca le risorse sandbox create per il controllo.
\item
  BlueAgent risponde con stato \passthrough{\lstinline!HTTP 200 OK!}
  comunicando l\textquotesingle esito della verifica.
\end{enumerate}

Se si verifica un errore le risorse utilizzate vengono rilasciate e
viene comunicato l\textquotesingle errore.

\subsection{\texorpdfstring{\texttt{UC-04}: Deploy di un unico container
per il team \(k\) {[}DA
RIVEDERE{]}}{UC-04: Deploy di un unico container per il team k {[}DA RIVEDERE{]}}}\label{uc-04-deploy-di-un-unico-container-per-il-team-k-da-rivedere}

\begin{itemize}
\tightlist
\item
  \textbf{Attore}: CTF Manager
\item
  \textbf{Obiettivo}: Avviare un singolo container specifico per il team
  \(k\), a partire dalla definizione della challenge che lo contiene.
\item
  \textbf{Precondizioni}: \passthrough{\lstinline!UC-01!}.
  E\textquotesingle{} stato eseguito con successo il deploy
  dell\textquotesingle infrastruttura di rete per il team \(k\).
  L\textquotesingle immagine del container richiesta e\textquotesingle{}
  presente nel catalogo in stato \passthrough{\lstinline!ready!}.
\end{itemize}

Flusso delle azioni:

\begin{enumerate}
\def\labelenumi{\arabic{enumi}.}
\tightlist
\item
  Il CTF Manager invia la richiesta
  \passthrough{\lstinline!POST /teams/\{team\_k\}/machines!} indicando
  l\textquotesingle ID del team nel path e passando la configurazione
  del container.
\item
  BlueAgent verifica che il team \(k\) esista e che il bridge
  \passthrough{\lstinline!br-team-k!} sia attivo.
\item
  BlueAgent verifica che l\textquotesingle immagine richiesta sia
  disponibile nel catalogo e pronta (\passthrough{\lstinline!ready!}).
\item
  BlueAgent avvia il container.
\item
  BlueAgent risponde con stato
  \passthrough{\lstinline!HTTP 201 Created!} confermando il deploy del
  singolo container.
\end{enumerate}

In caso di fallimento durante l\textquotesingle avvio, BlueAgent esegue
il rollback eliminando l\textquotesingle istanza creata e risponde con
un errore HTTP.

\subsection{\texorpdfstring{\texttt{UC-05}: Dismissione e rilascio di un
unico container per il team \(k\) {[}DA
RIVEDERE{]}}{UC-05: Dismissione e rilascio di un unico container per il team k {[}DA RIVEDERE{]}}}\label{uc-05-dismissione-e-rilascio-di-un-unico-container-per-il-team-k-da-rivedere}

\begin{itemize}
\tightlist
\item
  \textbf{Attore}: CTF Manager
\item
  \textbf{Obiettivo}: Arrestare ed eliminare un singolo container attivo
  per il team \(k\).
\item
  \textbf{Precondizioni}: Il container specificato appartiene al team
  \(k\) ed e\textquotesingle{} presente nel sistema.
\end{itemize}

Flusso delle azioni:

\begin{enumerate}
\def\labelenumi{\arabic{enumi}.}
\tightlist
\item
  Il CTF Manager invia la richiesta
  \passthrough{\lstinline!DELETE /teams/\{team\_k\}/machines/\{name\}!}
  specificando il nome del container da eliminare.
\item
  BlueAgent individua l\textquotesingle istanza del container:

  \begin{enumerate}
  \def\labelenumii{\arabic{enumii}.}
  \tightlist
  \item
    Arresta l\textquotesingle istanza del container.
  \item
    Elimina l\textquotesingle istanza, scollegandola dal bridge
    \passthrough{\lstinline!br-team-k!} e rilasciando le risorse
    allocate.
  \item
    Rimuove il record associato al container dal database.
  \end{enumerate}
\item
  BlueAgent risponde con stato
  \passthrough{\lstinline!HTTP 204 No Content!}.
\end{enumerate}

In caso di errore durante la rimozione del container, BlueAgent risponde
con errore \passthrough{\lstinline!5XX!}.

\subsection{\texorpdfstring{\texttt{UC-06}: Dismissione e rilascio di
una challenge per il team
\(k\)}{UC-06: Dismissione e rilascio di una challenge per il team k}}\label{uc-06-dismissione-e-rilascio-di-una-challenge-per-il-team-k}

\begin{itemize}
\tightlist
\item
  \textbf{Attore}: CTF Manager
\item
  \textbf{Obiettivo}: Rimuovere tutti i container associati a una
  determinata challenge per il team \(k\).
\item
  \textbf{Precondizioni}: I container associati alla challenge
  specificata sono attivi per il team \(k\).
\end{itemize}

Flusso delle azioni:

\begin{enumerate}
\def\labelenumi{\arabic{enumi}.}
\tightlist
\item
  Il CTF Manager invia la richiesta
  \passthrough{\lstinline!DELETE /teams/\{team\_k\}/machines!}
  specificando nel payload la configurazione della challenge.
\item
  Per ciascun container specificato nella configurazione della
  challenge:

  \begin{enumerate}
  \def\labelenumii{\arabic{enumii}.}
  \tightlist
  \item
    Arresta l\textquotesingle istanza.
  \item
    Elimina l\textquotesingle istanza.
  \item
    Cancella il record del container dal database.
  \end{enumerate}
\item
  BlueAgent risponde con stato
  \passthrough{\lstinline!HTTP 204 No Content!}.
\end{enumerate}

In caso di problemi durante il teardown di uno o piu container,
BlueAgent restituisce un errore HTTP con i dettagli
dell\textquotesingle anomalia.

\subsection{\texorpdfstring{\texttt{UC-07}: Dismissione e rilascio di un
intero team
\(k\)}{UC-07: Dismissione e rilascio di un intero team k}}\label{uc-07-dismissione-e-rilascio-di-un-intero-team-k}

\begin{itemize}
\tightlist
\item
  \textbf{Attore}: CTF Manager
\item
  \textbf{Obiettivo}: Rimuovere completamente
  l\textquotesingle infrastruttura e tutte le risorse associate al team
  \(k\).
\item
  \textbf{Precondizioni}: Il team \(k\) e\textquotesingle{} registrato
  nel sistema.
\end{itemize}

Flusso delle azioni:

\begin{enumerate}
\def\labelenumi{\arabic{enumi}.}
\tightlist
\item
  Il CTF Manager invia la richiesta
  \passthrough{\lstinline!DELETE /teams/\{team\_k\}!} specificando
  l\textquotesingle identificativo \passthrough{\lstinline!team\_k!}.
\item
  BlueAgent esegue la procedura di dismissione completa:

  \begin{enumerate}
  \def\labelenumii{\arabic{enumii}.}
  \tightlist
  \item
    Arresta ed elimina tutti i container ancora attivi per il team
    \(k\).
  \item
    Esegue il teardown dell\textquotesingle interfaccia WireGuard
    \passthrough{\lstinline!wg-k!} e le chiavi crittografiche associate
    ai Peer.
  \item
    Elimina le regole di routing e rimuove il bridge
    \passthrough{\lstinline!br-team-k!}.
  \item
    Cancella dal database il record del team \(k\) e tutti i record dei
    relativi membri e peer WireGuard.
  \end{enumerate}
\item
  BlueAgent risponde con stato
  \passthrough{\lstinline!HTTP 204 No Content!}.
\end{enumerate}

Se il team non esiste o e\textquotesingle{} gia stato dismesso,
l\textquotesingle operazione completa con successo senza sollevare
errori non gestiti.

\subsection{\texorpdfstring{\texttt{UC-08}: Compilazione e Import di
un\textquotesingle immagine per
container}{UC-08: Compilazione e Import di un\textquotesingle immagine per container}}\label{uc-08-compilazione-e-import-di-unimmagine-per-container}

\begin{itemize}
\tightlist
\item
  \textbf{Attore}: CTF Manager
\item
  \textbf{Obiettivo}: Inviare la specifica di
  un\textquotesingle immagine in formato YAML ed avviare la compilazione
  e l\textquotesingle importazione nel catalogo delle immagini.
\item
  \textbf{Precondizione}: BlueAgent e\textquotesingle{} attivo.
\end{itemize}

Flusso delle azioni:

\begin{enumerate}
\def\labelenumi{\arabic{enumi}.}
\tightlist
\item
  Il CTF Manager invia una richiesta
  \passthrough{\lstinline!POST /images!} inviando una form multipart con
  il file di definizione YAML.
\item
  BlueAgent valida la specifica YAML e ne calcola l\textquotesingle hash
  crittografico.
\item
  BlueAgent verifica lo stato dell\textquotesingle immagine nel
  database:

  \begin{enumerate}
  \def\labelenumii{\arabic{enumii}.}
  \tightlist
  \item
    Se l\textquotesingle immagine con lo stesso alias o hash esiste gia,
    risponde immediatamente con \passthrough{\lstinline!HTTP 200 OK!}.
  \item
    Altrimenti, registra un nuovo record nel database.
  \end{enumerate}
\item
  BlueAgent avvia la compilazione tramite un worker asincrono:

  \begin{enumerate}
  \def\labelenumii{\arabic{enumii}.}
  \tightlist
  \item
    Importa l\textquotesingle immagine generata nel catalogo dei
    container.
  \item
    Assegna l\textquotesingle alias all\textquotesingle immagine
    importata e aggiorna il record nel database.
  \end{enumerate}
\item
  BlueAgent restituisce immediatamente al CTF Manager lo stato
  \passthrough{\lstinline!HTTP 202 Accepted!}.
\end{enumerate}

\subsection{\texorpdfstring{\texttt{UC-09}: Dismissione di
un\textquotesingle immagine per container {[}DA
RIVEDERE{]}}{UC-09: Dismissione di un\textquotesingle immagine per container {[}DA RIVEDERE{]}}}\label{uc-09-dismissione-di-unimmagine-per-container-da-rivedere}

\begin{itemize}
\tightlist
\item
  \textbf{Attore}: CTF Manager
\item
  \textbf{Obiettivo}: Rimuovere un\textquotesingle immagine non piu
  necessaria dal catalogo di BlueAgent.
\item
  \textbf{Precondizioni}: L\textquotesingle immagine identificata
  dall\textquotesingle alias esiste nel sistema e non
  e\textquotesingle{} attualmente utilizzata da alcun container attivo.
\end{itemize}

Flusso delle azioni:

\begin{enumerate}
\def\labelenumi{\arabic{enumi}.}
\tightlist
\item
  Il CTF Manager invia la richiesta
  \passthrough{\lstinline!DELETE /images/\{alias\}!} specificando
  l\textquotesingle alias dell\textquotesingle immagine da eliminare.
\item
  BlueAgent verifica lo stato e i vincoli nel database:

  \begin{enumerate}
  \def\labelenumii{\arabic{enumii}.}
  \tightlist
  \item
    Verifica che l\textquotesingle immagine sia registrata a catalogo
    (altrimenti risponde con
    \passthrough{\lstinline!HTTP 404 Not Found!}).
  \item
    Verifica se vi siano container attivi che referenziano
    l\textquotesingle immagine. Se l\textquotesingle immagine
    e\textquotesingle{} in uso da container attivi, risponde con
    \passthrough{\lstinline!HTTP 409 Conflict!}.
  \end{enumerate}
\item
  Se l\textquotesingle immagine non e\textquotesingle{} utilizzata:

  \begin{enumerate}
  \def\labelenumii{\arabic{enumii}.}
  \tightlist
  \item
    Elimina l\textquotesingle immagine dal runtime dei container.
  \item
    Elimina il record dell\textquotesingle immagine dal database.
  \end{enumerate}
\item
  BlueAgent risponde con stato
  \passthrough{\lstinline!HTTP 204 No Content!}.
\end{enumerate}

\subsection{\texorpdfstring{\texttt{UC-10}: Ispezione dei Peer
WireGuard}{UC-10: Ispezione dei Peer WireGuard}}\label{uc-10-ispezione-dei-peer-wireguard}

\begin{itemize}
\tightlist
\item
  \textbf{Attore}: CTF Manager
\item
  \textbf{Obiettivo}: Monitorare lo stato della connettività VPN, le
  statistiche di traffico (byte ricevuti e trasmessi), gli handshake e
  lo stato online dei membri/giocatori del team \(k\).
\item
  \textbf{Precondizioni}: Il team \(k\) e\textquotesingle{} registrato e
  la relativa interfaccia WireGuard \passthrough{\lstinline!wg-k!}
  e\textquotesingle{} attiva sul sistema.
\end{itemize}

Flusso delle azioni:

\begin{enumerate}
\def\labelenumi{\arabic{enumi}.}
\tightlist
\item
  Il CTF Manager invia la richiesta
  \passthrough{\lstinline!GET /teams/\{team\_k\}/peers!} specificando
  l\textquotesingle id \(k\) del team.
\item
  BlueAgent interroga l\textquotesingle interfaccia di rete per estrarre
  le statistiche live dei peer.
\item
  BlueAgent interroga il database per recuperare i record dei membri
  registrati per il team \(k\).
\item
  BlueAgent risponde con stato \passthrough{\lstinline!HTTP 200 OK!}
  restituendo le statistiche dei peer raccolte.
\end{enumerate}

\subsection{\texorpdfstring{\texttt{UC-11}: Ispezione delle macchine in
esecuzione}{UC-11: Ispezione delle macchine in esecuzione}}\label{uc-11-ispezione-delle-macchine-in-esecuzione}

\begin{itemize}
\tightlist
\item
  \textbf{Attore}: CTF Manager
\item
  \textbf{Obiettivo}: Ottenere l\textquotesingle elenco e lo stato di
  tutti i container attualmente istanziati per il team \(k\).
\item
  \textbf{Precondizioni}: Il team \(k\) e\textquotesingle{} registrato
  nel sistema.
\end{itemize}

Flusso delle azioni:

\begin{enumerate}
\def\labelenumi{\arabic{enumi}.}
\tightlist
\item
  Il CTF Manager invia la richiesta
  \passthrough{\lstinline!GET /teams/\{team\_k\}/machines!} specificando
  l\textquotesingle id \(k\) del team.
\item
  BlueAgent interroga il database per recuperare i container registrati
  per il team \(k\).
\item
  BlueAgent verifica lo stato effettivo di ciascuna istanza.
\item
  BlueAgent risponde con stato \passthrough{\lstinline!HTTP 200 OK!}
  restituendo la lista dei container attivi.
\end{enumerate}

\subsection{\texorpdfstring{\texttt{UC-12}: Registrazione e
configurazione WireGuard per un membro del team
\(k\)}{UC-12: Registrazione e configurazione WireGuard per un membro del team k}}\label{uc-12-registrazione-e-configurazione-wireguard-per-un-membro-del-team-k}

\begin{itemize}
\tightlist
\item
  \textbf{Attore}: CTF Manager
\item
  \textbf{Obiettivo}: Registrare un giocatore come membro del team
  \(k\), allocare un indirizzo IP nella subnet VPN, generare le chiavi
  crittografiche WireGuard e restituire il file di configurazione
  \passthrough{\lstinline!.conf!}.
\item
  \textbf{Precondizioni}: BlueAgent e\textquotesingle{} attivo. Se
  l\textquotesingle infrastruttura di rete del team \(k\) non
  e\textquotesingle{} ancora attiva, viene allestita contestualmente.
\end{itemize}

Flusso delle azioni:

\begin{enumerate}
\def\labelenumi{\arabic{enumi}.}
\tightlist
\item
  Il CTF Manager invia la richiesta
  \passthrough{\lstinline!POST /teams/\{team\_k\}/members!} fornendo
  nome ed id dell\textquotesingle utente.
\item
  BlueAgent verifica se l\textquotesingle infrastruttura del team \(k\)
  e\textquotesingle{} attiva; se assente, esegue
  \passthrough{\lstinline!UC-01!}.
\item
  BlueAgent verifica se il membro e\textquotesingle{} gia registrato nel
  database per il team \(k\):

  \begin{enumerate}
  \def\labelenumii{\arabic{enumii}.}
  \tightlist
  \item
    Se il membro esiste gia, BlueAgent fornisce la configurazione
    rispondendo con \passthrough{\lstinline!HTTP 200 OK!}.
  \item
    Se il membro e\textquotesingle{} nuovo:

    \begin{enumerate}
    \def\labelenumiii{\arabic{enumiii}.}
    \tightlist
    \item
      Alloca un indirizzo IP VPN libero per il team. Se lo spazio degli
      indirizzi e\textquotesingle{} esaurito, risponde segnalando
      l\textquotesingle esaurimento degli indirizzi.
    \item
      Se non fornita, genera una nuova coppia di chiavi asimmetriche.
    \item
      Registra il nuovo peer sull\textquotesingle interfaccia
      \passthrough{\lstinline!wg-k!}.
    \item
      Salva il record del membro nel database.
    \end{enumerate}
  \end{enumerate}
\item
  BlueAgent risponde con stato
  \passthrough{\lstinline!HTTP 201 Created!} restituendo la
  configurazione WireGuard generata.
\end{enumerate}

\subsection{\texorpdfstring{\texttt{UC-13}: Revoca di un membro
WireGuard per il team
\(k\)}{UC-13: Revoca di un membro WireGuard per il team k}}\label{uc-13-revoca-di-un-membro-wireguard-per-il-team-k}

\begin{itemize}
\tightlist
\item
  \textbf{Attore}: CTF Manager
\item
  \textbf{Obiettivo}: Rimuovere un membro registrato dal team \(k\),
  eliminando il relativo peer dall\textquotesingle interfaccia WireGuard
  \passthrough{\lstinline!wg-k!} e liberando l\textquotesingle indirizzo
  IP VPN allocato per futuri riutilizzi.
\item
  \textbf{Precondizioni}: Il membro specificato e\textquotesingle{}
  registrato nel team \(k\).
\end{itemize}

Flusso delle azioni:

\begin{enumerate}
\def\labelenumi{\arabic{enumi}.}
\tightlist
\item
  Il CTF Manager invia la richiesta
  \passthrough{\lstinline!DELETE /teams/\{team\_k\}/members/\{member\_id\}!}
  specificando l\textquotesingle identificativo del team e del membro.
\item
  BlueAgent individua il membro nel database:

  \begin{enumerate}
  \def\labelenumii{\arabic{enumii}.}
  \tightlist
  \item
    Se non trovato, risponde con
    \passthrough{\lstinline!HTTP 404 Not Found!}.
  \end{enumerate}
\item
  BlueAgent rimuove il peer dall\textquotesingle interfaccia
  \passthrough{\lstinline!wg-k!}.
\item
  BlueAgent elimina il record del membro dal database.
\item
  BlueAgent risponde con stato
  \passthrough{\lstinline!HTTP 204 No Content!}.
\end{enumerate}
